\begin{frame}\frametitle{What is a security goal?}
\pause
A \emph{security goal} is a goal that we assign for a given protocol to evaluate if it is sufficiently secure for application.

\pause
Security goals formalize our notion of ``secure''.

\bigskip
\pause
Not every protocol has the same requirement for security.

Common Security Goals
\begin{itemize}
\pause
\item \textbf{Confidentiality:} In a two-party communication protocol, confidentiality holds if no third party is able to access the information sent.

\pause
\item \textbf{Agreement:} In a two-party communication protocol, agreement holds if both parties agree on some set of data values in the protocol.
These values are then used to derive the key.
\end{itemize}

\end{frame}

\begin{frame}\frametitle{Other common security goals}
\begin{enumerate}
\item Aliveness
\item Weak Agreement
\item Non-injective Agreement
\item A specific relay party cannot read the message
\item The message must be timestamped
\item Perfect forward secrecy
\item Crypto primitives are resistant to a predefined attack
\end{enumerate}
\end{frame}